% Source-PDF: 数学 MATHEMATICS/线性规划与单纯形算法 - 吴一凡.pdf
\documentclass[11pt]{article}
\usepackage{oipl-vn}

\newcommand{\Oh}{\mathcal{O}}
\newcommand{\ST}{\text{S.T.}}
\newcommand{\INF}{\mathrm{INF}}

\title{Quy hoạch tuyến tính và thuật toán đơn hình}
\author{Wu Yifan\\Khối trung học phổ thông, Trường Phổ thông trực thuộc Đại học Cát Lâm}
\date{3 tháng 3, 2015}

\begin{document}
\maketitle

\origtitle{Linear Programming and Simplex Algorithm}
\sourcepath{Mathematics / Linear programming and simplex algorithm - Wu Yifan}

\begin{abstract}
Bài giảng nhắc lại dạng chuẩn của quy hoạch tuyến tính, cách đưa một mô
hình tuyến tính bất kỳ về dạng chuẩn và dạng lỏng, rồi trình bày thuật toán
đơn hình qua khái niệm biến cơ sở, biến phi cơ sở, nghiệm khả thi, phép
xoay trục và quy tắc Bland. Phần cuối nêu một số ứng dụng trong bài toán thi
lập trình, đặc biệt là mối liên hệ với mô hình luồng mạng, đối ngẫu tuyến
tính và ví dụ BZOJ3876.
\end{abstract}

\tableofcontents

\section{Giới thiệu}

Ở bậc phổ thông, nhiều người đã gặp quy hoạch tuyến tính và biết đại khái
nó là gì: có một số biến, cần tìm giá trị tốt nhất của một hàm theo các biến
đó dưới một số ràng buộc. Cả ràng buộc lẫn hàm mục tiêu đều tuyến tính.

Nếu các bài quy hoạch tuyến tính gặp trong chương trình phổ thông đều được
giải bằng hình học, thì cách đó trở nên rất khó chịu khi số chiều tăng lên.
Vì vậy ta dùng thuật toán đơn hình để xử lý trực tiếp hơn. Các chứng minh
chi tiết trong bài giảng gốc được giữ ở mức trực giác.

\section{Dạng chuẩn}

Ta định nghĩa một dạng chuẩn cho quy hoạch tuyến tính:
\[
  \begin{array}{ll}
  \text{Maximize} & c^T x,\\
  \ST & Ax\le b,\quad x\ge 0.
  \end{array}
\]
Ở đây giả sử có \(n\) biến và \(m\) ràng buộc. Vector \(x\) là vector cột
chứa các biến; \(c\) là vector hệ số của hàm mục tiêu; \(A\) là ma trận hệ
số của các ràng buộc; \(b\) là vector cột hằng số ở vế phải. Điều kiện
\(x\ge 0\) nghĩa là mọi thành phần của \(x\) đều không âm. Kích thước tương
ứng là
\[
  c^T[1,n],\qquad x[n,1],\qquad A[m,n],\qquad b[m,1].
\]

Tất nhiên không phải mọi bài toán quy hoạch tuyến tính ban đầu đều đã ở
dạng này. Ta lần lượt xét các phép biến đổi đưa mô hình tổng quát về dạng
chuẩn, và các phép biến đổi này không làm thay đổi giá trị tối ưu cần tìm.

\subsection{Hàm mục tiêu}

Nếu bài toán yêu cầu cực tiểu thay vì cực đại, ta đổi dấu toàn bộ hệ số của
hàm mục tiêu, giải bài toán cực đại mới, rồi đổi dấu đáp án. Nói cách khác,
\[
  \min c^T x = - \max (-c)^T x.
\]

\subsection{Chiều của bất đẳng thức}

Trong dạng chuẩn, mọi ràng buộc đều có dạng \(\le\). Nếu ràng buộc có dạng
\(\ge\), ta nhân cả hai vế với \(-1\) để đổi về \(\le\). Nếu ràng buộc là
đẳng thức, chẳng hạn \(a=b\), ta thay bằng hai ràng buộc
\[
  a\le b,\qquad b\le a.
\]

\subsection{Ràng buộc trên từng biến}

Mỗi biến trong dạng chuẩn phải không âm. Với các ràng buộc riêng trên một
biến, ta thay biến như sau.

Nếu \(x\le a\), đặt
\[
  x'=a-x,\qquad x'\ge 0.
\]
Nếu \(x\ge a\), đặt
\[
  x'=x-a,\qquad x'\ge 0.
\]
Nếu \(x\) không bị chặn dấu, đặt
\[
  x=x'-x'',\qquad x'\ge 0,\quad x''\ge 0.
\]
Sau các phép thay thế trên, mọi quy hoạch tuyến tính đều có thể chuyển về
dạng chuẩn.

\section{Dạng lỏng}

Để áp dụng đơn hình thuận tiện hơn, ta đổi các bất đẳng thức trong dạng
chuẩn thành đẳng thức bằng cách thêm biến phụ. Với ràng buộc thứ \(i\):
\[
  \sum_{j=1}^{n} A_{i,j}x_j \le b_i,
\]
thêm biến \(x_{n+i}\) và viết
\[
  x_{n+i}+\sum_{j=1}^{n} A_{i,j}x_j=b_i,\qquad x_{n+i}\ge 0.
\]
Như vậy mọi ràng buộc đã trở thành đẳng thức. Tổng cộng có \(n+m\) biến và
mọi biến vẫn không âm. Thuật toán đơn hình sẽ được mô tả trên dạng lỏng này.

\section{Một số định nghĩa}

\term{Biến cơ sở} là biến được viết đầu tiên trong mỗi phương trình ràng
buộc. Trong dạng lỏng ban đầu, \(x_{n+1},\ldots,x_{n+m}\) là các biến cơ sở,
nên có đúng \(m\) biến cơ sở.

\term{Biến phi cơ sở} là các biến còn lại. Ban đầu
\(x_1,\ldots,x_n\) là biến phi cơ sở, nên có đúng \(n\) biến phi cơ sở. Hàm
mục tiêu chỉ liên quan trực tiếp đến các biến phi cơ sở.

Ký hiệu tập biến cơ sở là \(B\), tập biến phi cơ sở là \(N\). Hai tập này
thay đổi trong quá trình thuật toán.

\term{Nghiệm khả thi} là một bộ giá trị làm thỏa mãn tất cả ràng buộc.
\term{Nghiệm tối ưu} là nghiệm khả thi làm hàm mục tiêu đạt giá trị tốt
nhất. Mục tiêu của quy hoạch tuyến tính là tìm một nghiệm tối ưu.

\section{Thuật toán đơn hình}

Trước hết giả sử nghiệm cơ sở hiện tại là khả thi, tức là khi mọi biến phi
cơ sở bằng \(0\), các biến cơ sở thu được đều không âm. Trường hợp không có
tính chất này sẽ được xét sau.

Nếu tồn tại một biến phi cơ sở có hệ số dương trong hàm mục tiêu, tăng biến
đó có thể làm giá trị mục tiêu tăng. Nếu mọi hệ số của biến phi cơ sở đều
không dương, nghiệm hiện tại đã tối ưu: khi các biến phi cơ sở không âm,
không thể tăng mục tiêu thêm. Lúc này giá trị tối ưu chính là hằng số còn
lại trong hàm mục tiêu.

Giả sử ta chọn biến phi cơ sở thứ \(e\), ký hiệu \(x_{N_e}\), để tăng.
Biến này không thể tăng vô hạn nếu một ràng buộc trở nên chặt trước. Xét
ràng buộc thứ \(i\):
\[
  x_{B_i}+\sum_{j=1}^{n} A_{i,j}x_{N_j}=b_i.
\]
Nếu \(A_{i,e}>0\), để giữ \(x_{B_i}\ge 0\) ta cần
\[
  x_{N_e}\le \frac{b_i}{A_{i,e}}.
\]
Ta chọn ràng buộc chặt nhất. Giả sử đó là ràng buộc thứ \(l\), tương ứng với
biến cơ sở \(x_{B_l}\). Nếu không có ràng buộc nào như vậy, tức giá trị nhỏ
nhất là \(\INF\), thì \(x_{N_e}\) có thể tăng vô hạn và bài toán không bị
chặn trên.

Vì ràng buộc thứ \(l\) là chặt nhất, đẳng thức tương ứng thật sự đạt dấu
bằng. Ta thực hiện phép xoay trục: giải phương trình đó theo \(x_{N_e}\),
rồi thay \(x_{N_e}\) vào mọi phương trình khác và vào hàm mục tiêu. Khi đó
\(x_{B_l}\) trở thành biến phi cơ sở thứ \(e\), còn \(x_{N_e}\) trở thành
biến cơ sở thứ \(l\). Nói cách khác, hai biến này đổi vai trò cho nhau, và
cả hai tập \(B,N\) đều thay đổi.

\subsection{Ví dụ xoay trục}

Xét quy hoạch tuyến tính:
\[
\begin{array}{ll}
\text{Maximize} & 3x_1+x_2+2x_3,\\
\ST & x_4+x_1+x_2+3x_3=30,\\
    & x_5+2x_1+2x_2+5x_3=24,\\
    & x_6+4x_1+x_2+2x_3=36.
\end{array}
\]
Ban đầu
\[
  B=\{4,5,6\},\qquad N=\{1,2,3\}.
\]
Ta chọn biến phi cơ sở \(x_1\) để tăng. Từ ba ràng buộc nhận được lần lượt
\[
  x_1\le 30,\qquad x_1\le 12,\qquad x_1\le 9.
\]
Ràng buộc thứ ba chặt nhất, nên ta đổi biến cơ sở thứ ba với biến phi cơ sở
thứ nhất. Từ ràng buộc thứ ba:
\[
  x_1=9-\frac{x_6}{4}-\frac{x_2}{4}-\frac{x_3}{2}.
\]
Thay biểu thức này vào những nơi còn lại, ta được quy hoạch tuyến tính mới:
\[
\begin{array}{ll}
\text{Maximize} & -\dfrac{3x_6}{4}+\dfrac{x_2}{4}
                  +\dfrac{x_3}{2}+27,\\[0.5em]
\ST & x_4-\dfrac{x_6}{4}+\dfrac{3x_2}{4}
        +\dfrac{5x_3}{2}=21,\\[0.5em]
    & x_5-\dfrac{x_6}{2}+\dfrac{3x_2}{2}+4x_3=6,\\[0.5em]
    & x_1+\dfrac{x_6}{4}+\dfrac{x_2}{4}+\dfrac{x_3}{2}=9.
\end{array}
\]
Lúc này
\[
  B=\{4,5,1\},\qquad N=\{6,2,3\}.
\]
Hằng số trong hàm mục tiêu đã tăng. Lặp quá trình này, thuật toán sẽ tìm
được một nghiệm tối ưu hoặc kết luận miền khả thi không bị chặn theo hướng
cực đại.

\subsection{Mã xoay trục}

Slide gốc chèn đoạn mã xoay trục dưới dạng ảnh chụp màn hình. Ở đây nội dung
được chép lại thành văn bản LaTeX để có thể đọc, tìm kiếm và biên dịch ổn
định; tên biến được giữ theo slide:
\begin{verbatim}
inline void pivot(int l,int e){
    int i,j,k;
    b[l]/=A[l][e];
    for(i=1;i<=n;++i)if(i!=e)A[l][i]/=A[l][e];
    A[l][e]=1/A[l][e];

    for(i=1;i<=m;++i)if(i!=l&&fabs(A[i][e])>eps){
        b[i]-=A[i][e]*b[l];
        for(j=1;j<=n;++j)if(j!=e)A[i][j]-=A[i][e]*A[l][j];
        A[i][e]-=A[i][e]*A[l][e];
    }

    v+=c[e]*b[l];
    for(i=1;i<=n;++i)if(i!=e)c[i]-=c[e]*A[l][i];
    c[e]=-c[e]*A[l][e];

    swap(B[l],N[e]);
}
\end{verbatim}

Cần chú ý quy tắc Bland: khi chọn biến phi cơ sở đi vào cơ sở, hãy chọn biến
có hệ số dương trong mục tiêu và có chỉ số nhỏ nhất. Cách chọn này tránh các
vòng lặp do trường hợp suy biến.

Mỗi lần xoay trục có độ phức tạp \(\Oh(nm)\). Thuật toán đơn hình không phải
thuật toán thời gian đa thức theo nghĩa chặt, và vẫn có thể bị dữ liệu đặc
biệt làm chậm đến mức mũ. Tuy vậy trong thực hành lập trình thi đấu, nếu dữ
liệu không cố ý khai thác điểm yếu này, nó thường chạy đủ tốt.

\section{Ứng dụng}

Thuật toán đơn hình dùng để làm gì? Câu trả lời thực dụng trong slide là:
giải bài tập. Các slide 21--30 nhấn mạnh ý này bằng một chuỗi overlay lặp
ý ``làm bài áp dụng trực tiếp''. Tuy nhiên trước các bài quy hoạch tuyến tính
mạnh, những bài chỉ yêu cầu áp dụng trực tiếp không nhiều.

Một nguồn mô hình quan trọng là luồng mạng. Các bài như luồng cực đại,
luồng chi phí nhỏ nhất, luồng có chặn trên và chặn dưới đều có thể xem là
quy hoạch tuyến tính. Do giới hạn thời gian, slide chỉ nhắc hướng này và
gợi ý tham khảo tài liệu ``Quy hoạch tuyến tính và luồng mạng'' của Cao
Qinxiang.

\section{Ví dụ BZOJ3876}

Xét bài BZOJ3876, Ahoi2014, tên thường dịch là ``Cốt truyện nhánh''. Nếu bỏ
qua mô hình luồng mạng, ta có thể dùng quy hoạch tuyến tính trực tiếp.

Gọi \(x_i\) là số lần cạnh thứ \(i\) được đi qua. Thêm một siêu nguồn nối
vào đỉnh \(1\), và mỗi đỉnh nối đến siêu đích. Với mọi đỉnh trung gian, số
lần đi vào bằng số lần đi ra. Ngoài ra \(x_i\ge 1\) đối với các cạnh gốc của
bài toán; điều kiện này không áp dụng cho các cạnh thêm vào.

Ta có thể đặt từng ràng buộc riêng vào hệ lớn. Hàm mục tiêu cũng rất rõ
ràng: cần cực tiểu tổng chi phí. Đổi dấu hệ số để biến thành bài cực đại.
Vấn đề là một số ràng buộc có dạng \(\ge 1\); sau khi đổi dấu chúng thành
\(\le -1\), nghiệm cơ sở ban đầu không còn khả thi.

\subsection{Đối ngẫu tuyến tính}

Với bài toán
\[
  \begin{array}{ll}
  \text{Maximize} & c^T x,\\
  \ST & Ax\le b,\quad x\ge 0,
  \end{array}
\]
bài toán đối ngẫu là
\[
  \begin{array}{ll}
  \text{Minimize} & b^T y,\\
  \ST & A^T y\ge c,\quad y\ge 0.
  \end{array}
\]
Hai bài toán có cùng giá trị tối ưu. Vì vậy ta có thể đối ngẫu trực tiếp
mô hình ban đầu để thu được một bài toán cực đại có nghiệm cơ sở ban đầu khả
thi, khỏi phải tự xây dựng nghiệm khả thi ban đầu.

Cách làm này vẫn có thể quá chậm. Trong mô hình trực tiếp, số biến và số
ràng buộc đều là \(\Oh(n+m)\); với bài này thời gian chạy được nhận xét là
khoảng ba lần ZKW chi phí nhỏ nhất, nên không thật sự phù hợp.

Điểm cần sửa là không đặt quá nhiều ràng buộc đơn lẻ cho từng phần tử. Ta có
thể gộp chúng vào các ràng buộc tổng, dùng biến thay thế dạng \(x'=x-1\),
để số ràng buộc của bài gốc còn \(\Oh(n)\). Sau khi đối ngẫu, hệ có
\(\Oh(n)\) biến và \(\Oh(n+m)\) ràng buộc, gần với mức có thể chấp nhận.
Nhược điểm là ở bước này có thể lại mất nghiệm khả thi ban đầu, nên cần xử
lý cẩn thận theo từng mô hình.

\section{Kết luận}

Quy hoạch tuyến tính có nhiều ứng dụng. Thuật toán đơn hình có cài đặt đơn
giản, độ phức tạp thực tế thường thấp và không dễ bị chặn bằng dữ liệu thông
thường; vì vậy nó là một thuật toán đáng cân nhắc trong thi lập trình.

Các bài gợi ý:
\[
  \text{BZOJ1061},\quad
  \text{BZOJ3112},\quad
  \text{BZOJ3265},\quad
  \text{BZOJ3876}.
\]

\end{document}
